\newpage
\chapter{CARP: Adaptive Range Partitioning for In-Situ Data Management}

\section{The In-Situ Data Management Landscape}

High-performance computing applications generate data at unprecedented scales, creating a
fundamental challenge in data management that has traditionally been addressed through a "persist
first, analyze later" paradigm.
As simulations scale to hundreds of thousands of cores and produce terabytes of output per
timestep, the costs of data movement have become prohibitive.
This has driven the emergence of in-situ analytics as a critical capability for large-scale
scientific computing.

The appeal of in-situ processing lies in its ability to leverage idle computational resources
during I/O phases to perform data analysis and organization tasks that would otherwise require
expensive post-processing.
Early demonstrations of this approach have shown significant promise, with systems like DeltaFS
successfully implementing hash-based partitioning at scales up to 131,072 cores.
However, the full potential of in-situ data management has been constrained by the limited
repertoire of available partitioning primitives.

\section{The Partitioning Primitive Gap}

Database systems have long recognized partitioning as a fundamental technique for organizing data
to enable efficient query processing.
Hash partitioning provides uniform distribution but offers limited query optimization benefits,
while range partitioning enables powerful query acceleration through partition pruning.
In the context of HPC applications, this distinction becomes critical when supporting complex
analytical workloads that require spatial or temporal locality.

Prior to this work, the in-situ data management toolkit was incomplete.
The DeltaFS project had demonstrated the viability of hash partitioning through its 3-hop shuffler,
an embedded indexing pipeline that used idle compute and network resources to partition data en
route to storage.
However, range partitioning remained elusive due to the dynamic nature of simulation data, where
the distribution of values evolves unpredictably over time.

The absence of viable in-situ range partitioning represented a significant gap in the toolkit.
Without this capability, applications were forced to choose between suboptimal hash partitioning
for in-situ processing or expensive post-processing for range-based organization.
This limitation prevented the full realization of query-efficient data layouts that could
accelerate downstream analytics.

\section{CARP Design and Architecture}

CARP addresses the range partitioning challenge through an adaptive control loop that continuously
monitors data distribution and adjusts partitioning decisions in real-time.
The core insight is that effective range partitioning requires not just knowledge of current data
characteristics, but the ability to adapt to changing conditions throughout the application's
execution.

The system implements an OODA (Observe-Orient-Decide-Act) loop specifically designed for the
partitioning problem.
The observation phase continuously samples incoming data to characterize its distribution.
The orientation phase analyzes this information in the context of current partitioning state and
query requirements.
The decision phase determines optimal partition boundaries, and the action phase implements these
decisions through a distributed renegotiation protocol.

\subsection{Adaptive Control Loop}

The adaptive control loop represents the heart of CARP's contribution to in-situ data management.
Unlike static partitioning schemes that rely on a priori knowledge of data distribution, CARP's
control loop enables dynamic adaptation to evolving conditions.
The system maintains lightweight statistics about data distribution and uses these to continuously
refine partitioning decisions.

The control loop operates on multiple timescales, enabling both fine-grained adjustments within
individual timesteps and coarse-grained adaptations across the application's execution.
This multi-scale approach ensures that the system can respond to both transient fluctuations and
long-term trends in data characteristics.

\subsection{Renegotiation Protocol}

The renegotiation protocol enables CARP to maintain consistent partitioning state across the
distributed system while adapting to changing conditions.
The protocol coordinates partition boundary updates across all participating nodes, ensuring that
the system maintains a coherent view of data organization despite ongoing adaptations.

The protocol is designed to minimize disruption to the ongoing computation while ensuring that
partitioning decisions remain globally consistent.
This is achieved through a combination of versioned partition boundaries and careful coordination
of state transitions across the distributed system.

\section{Implementation and Integration}

CARP's implementation demonstrates the practical viability of adaptive range partitioning in real
HPC environments.
The system is designed to integrate seamlessly with existing simulation workflows, requiring
minimal modifications to application code while providing significant improvements in data
organization capabilities.

The implementation addresses several key challenges in deploying adaptive partitioning at scale.
Resource management ensures that the control loop operations do not interfere with the primary
computation, while the integration framework provides clean interfaces for incorporating CARP into
diverse application domains.

\section{Experimental Validation}

The experimental validation of CARP demonstrates its effectiveness across multiple dimensions.
Performance evaluations show that the system can achieve query acceleration benefits comparable to
offline range partitioning while maintaining minimal runtime overhead.
Scalability studies confirm that the adaptive control loop can operate effectively at large scales
without introducing prohibitive coordination costs.

The results establish that in-situ range partitioning is not only viable but provides substantial
benefits over alternative approaches.
Query performance improvements of up to an order of magnitude are achieved for range-based
analytical workloads, while runtime overhead remains below 5

\section{Lessons
  and Implications}

CARP's success in addressing the range partitioning challenge
has broader implications for in-situ data management and HPC systems design.
% for typical simulation workloads.
The work demonstrates that adaptive control loops can be successfully applied to complex
distributed system challenges, providing a foundation for more sophisticated in-situ processing
capabilities.

The completion of the partitioning primitive toolkit through CARP's contribution enables a new
class of in-situ analytics workflows that can achieve query-efficient data organization without
expensive post-processing.
This capability represents a significant step toward realizing the full potential of in-situ data
management in HPC applications.

Perhaps most importantly, CARP establishes the viability of real-time observation and control as a
fundamental pattern for in-situ processing.
The adaptive control loop demonstrated in the data management domain provides a template that can
be applied to other HPC challenges, setting the stage for the unified framework approach explored
in subsequent chapters.

